\section{Conclusiones de la revisión sistemática de la literatura}

La presente revisión sistemática de la literatura abordó de manera rigurosa la identificación, selección y análisis de la evidencia científica relacionada con el dominio de estudio definido. Mediante la aplicación de una estrategia de búsqueda estructurada, complementada con técnicas de \textit{snowballing} y criterios de inclusión y exclusión previamente establecidos, fue posible conformar un conjunto final de 22 estudios relevantes para la investigación.

El análisis de los estudios seleccionados caracterizó el estado actual de la literatura, trazó las principales líneas de investigación abordadas por la comunidad científica y determinó el grado de cobertura de las preguntas de investigación planteadas. Los resultados obtenidos evidencian correspondencia entre las preguntas de investigación planteadas y la literatura recuperada, ya que la totalidad de los estudios contribuyó a responder la \textbf{RQ1}, mientras que la gran mayoría aportó evidencia para la \textbf{RQ2}.

Asimismo, el proceso de revisión consolidó una visión estructurada del área de investigación, identificando tendencias, enfoques y características comunes presentes en los trabajos analizados. Los hallazgos obtenidos sirven como referencia para el desarrollo de las etapas posteriores de la investigación, al aportar evidencia sobre el estado del arte del área estudiada.
